
\documentclass[12pt]{article}
\usepackage[spanish]{babel}
\usepackage[utf8]{inputenc}
\usepackage{amsmath}
\usepackage{amsfonts}
\usepackage{amssymb}
\usepackage{geometry}
\geometry{a4paper, margin=2.5cm}
\title{Ecuaciones Básicas de la Dinámica de Fluidos}
\author{Traducción técnica desde ``Theory and Numerics for Problems of Fluid Dynamics''}
\date{}

\begin{document}
\maketitle

\section*{1.1.6 Propiedades de los coeficientes de viscosidad}

Aquí $\mu$ y $\lambda$ son llamados el primer y segundo coeficiente de viscosidad, respectivamente. $\mu$ también se denomina viscosidad dinámica. En la teoría cinética de gases se derivan las condiciones

\[
\mu \geq 0, \quad 3\lambda + 2\mu \geq 0.
\]

Para gases monoatómicos, se cumple $3\lambda + 2\mu = 0$. Esta condición se utiliza usualmente incluso en el caso de gases más complejos.

\subsection*{1.1.7 La ecuación de energía}

La ley de conservación de la energía se expresa mediante la ecuación de energía. Para ello se define la energía total

\[
E = \rho\left(e + \frac{|v|^2}{2}\right),
\]

donde $e$ es la energía interna específica. La ecuación de energía tiene la forma

\[
\frac{\partial E}{\partial t} + \nabla \cdot (E \, v) = \rho f \cdot v + \nabla \cdot (T v) + \rho q - \nabla \cdot q.
\]

Para un fluido newtoniano se tiene

\[
\frac{\partial E}{\partial t} + \nabla \cdot (E \, v) = \rho f \cdot v - \nabla \cdot (p v) + \nabla \cdot (\lambda v \, \nabla \cdot v) + \nabla \cdot (2\mu D(v)v) + \rho q - \nabla \cdot q,
\]

donde $q$ es la densidad de fuentes térmicas y $q$ el flujo de calor, el cual depende de la temperatura según la ley de Fourier:

\[
q = -k \nabla \theta,
\]

donde $k \geq 0$ es el coeficiente de conductividad térmica.

\section*{1.2 Relaciones termodinámicas}

Para completar el sistema de leyes de conservación, deben incluirse ecuaciones adicionales derivadas de la termodinámica.

La temperatura absoluta $\theta$, la densidad $\rho$ y la presión $p$ se denominan variables de estado. Todas estas cantidades son funciones positivas. El gas se caracteriza por la ecuación de estado

\[
p = p(\rho, \theta),
\]

y la relación

\[
e = e(\rho, \theta).
\]

Consideraremos el llamado gas perfecto (o gas ideal), cuyas variables de estado satisfacen la ecuación de estado

\[
p = R \theta \rho,
\]

donde $R > 0$ es la constante del gas, la cual puede expresarse como

\[
R = c_p - c_v,
\]

donde $c_p$ y $c_v$ son los calores específicos a presión y volumen constantes, respectivamente. Los experimentos muestran que $c_p > c_v$, por lo que $R > 0$ y $c_p$, $c_v$ son constantes en un amplio rango de temperaturas. La cantidad

\[
\gamma = \frac{c_p}{c_v} > 1
\]

se denomina constante adiabática de Poisson. Por ejemplo, para el aire, $\gamma = 1.4$. La energía interna de un gas perfecto está dada por

\[
e = c_v \theta.
\]

\subsection*{1.2.1 Entropía}

Una de las cantidades termodinámicas importantes es la entropía $S$, definida por la relación

\[
\theta \, dS = de + p \, dV,
\]

donde $V = \frac{1}{\rho}$ es el llamado volumen específico. Esta identidad se deduce de la termodinámica bajo la suposición de que la energía interna es una función de $S$ y $V$: $e = e(S, V)$, lo que explica el significado de los diferenciales.

\textbf{Teorema 1.3.} Para un gas perfecto se tiene

\[
S = c_v \ln \left( \frac{p/p_0}{(\rho/\rho_0)^\gamma} \right) + \text{constante},
\]

donde $p_0$ y $\rho_0$ son valores de referencia fijos de presión y densidad.

\subsection*{1.2.2 Flujo barotrópico}

Se dice que el flujo es barotrópico si la presión puede expresarse como una función de la densidad:

\[
p = p(\rho),
\]

lo cual implica que $p(x,t) = p(\rho(x,t))$, o más brevemente, $p = p \circ \rho$. Suponemos que

\[
p : (0, +\infty) \to (0, +\infty),
\]

y que existe su derivada continua.

\end{document}
